\documentclass{article}
\usepackage[backend=biber]{biblatex}
\usepackage{graphicx}
\usepackage[colorlinks=true]{hyperref}
\usepackage{booktabs}
\usepackage{siunitx}
\usepackage[]{amsmath}
\usepackage{gensymb}
\usepackage{mathtools}
\usepackage{fancyref}
\addbibresource{/home/giorgio/Bibliography/bibliography.bib}
\hypersetup{
    colorlinks=true,
    linkcolor=blue,
    filecolor=magenta,  
    citecolor=blue,    
    urlcolor=cyan,
    pdftitle={Selective Laser Synthering for renewable polymers},
    bookmarks=true,
}

\author{Giorgio De Trane}
\title{\textbf{Selective Laser Synthering for renewable polymers}}

\begin{document}
    \setlength{\parindent}{0pt}
    \maketitle
    \begin{center}
        \includegraphics[width=0.8\textwidth]{Pictures/polito_logo.eps}\linebreak\newline
        \textit{Year 2021/2022}
    \end{center}

    \newpage
    \tableofcontents
    \newpage
    \section{Introduction\label{Intro}}

    \textit{Additive manufacturing} (AM) is a broad term that encompasses several manufacturing techniques, characterized by their additive nature, as the name suggests, 
    in contrast with more traditional subtractive processes.
    
    AM techniques are applied to a vast range of materials, including ceramics, polymers and metal alloys, some of which are specifically developed or optimized to 
    these kinds of applications. \\ 
    
    The main advantage of AM is the ability to produce complex shapes in a relatively short time. 
    These geometries are either too hard or even impossible to reproduce with subtractive manufacturing techniques, which often require multiple steps, using different 
    pieces of equipment, trained personnel, etc. \\

    Given the same material, the complex shapes allowed by AM can replace components made of multiple assembled parts with a unique solid piece of comparable or even better mechanical properties. 
    
    The inherent flexibility of AM often allows product designers to simplify or even entirely bypass the very strict CAD workflow (which is intrinsically tied to 
    traditional manufacturing processes) and make use of organic and/or generative modeling. 

    As a consequence of better design choices and minimal need of post-processing of AM objects, far less raw material is wasted, compared to subtractive manufacturing techniques, leading to 
    long term lowering of costs, faster design-to-market pipelines and, last but not least, lower emissions and environmental impact \autocite*{Recent_progress_polymers_AM}. 

    \subsection{Polymers and 3D printing}

    










    \clearpage

    \printbibliography
    
\end{document}